\section{The Autonomous Drift Problem}
\label{sec:rs-drift}

The Recursive Spawning Protocol is designed to prevent a specific structural failure: an agent that operates without any traceable authorization path to a human principal. This section provides the formal vocabulary for that failure, characterizes the conditions that produce it, catalogs its concrete manifestations, proves that depth limits alone cannot prevent it, and defines the formal model used throughout the remainder of the paper.

% ----------------------------------------------------------------------
\subsection{Formal Characterization}
\label{sec:rs-drift-formal}

We model a multi-agent system as a rooted directed acyclic graph (DAG) of agent nodes. Each node $n$ carries a parent pointer $\alpha(n)$ — the node that authorized its provisioning. The root of every authorized chain is a human principal $h \in \mathcal{H}$. A spawn chain $C = \langle n_0, n_1, \ldots, n_k \rangle$ satisfies $\alpha(n_i) = n_{i-1}$ for all $1 \leq i \leq k$ and $n_0 \in \mathcal{H}$. For notational convenience, we write $\alpha(n) = \bot$ when $n$ has no parent — this state applies only to root human principals.

\medskip
\noindent\textbf{Definition 1 (Orphaned Agent).}
An agent $n \in \mathcal{N}$ is \emph{orphaned} if its parent node is either terminated or unreachable:
\begin{enumerate}
  \item[(a)] \textbf{Terminated.} $\alpha(n)$ has exited operational state through normal completion, abnormal failure, or administrative shutdown.
  \item[(b)] \textbf{Unreachable.} $\alpha(n)$ does not respond to any defined communication protocol and this condition persists beyond the heartbeat timeout $\tau_{\mathrm{hb}}$.
\end{enumerate}
Orphaning is a structural property of the parent pointer — it does not depend on whether $n$ is aware of the parent's absence. An agent that is unaware it is orphaned is still orphaned.

\medskip
\noindent\textbf{Definition 2 (Drifted Agent).}
An agent $n$ is \emph{drifted} if the chain from $n$ to the root does not contain a human principal. Equivalently:
\begin{equation}
\mathrm{drifted}(n) \;\Longleftrightarrow\; \neg\exists\, h \in \mathcal{H} : \mathrm{chain-reaches}(n, h).
\label{eq:rs-drifted-def}
\end{equation}
where $\mathrm{chain-reaches}(n, h)$ holds when the transitive closure of $\alpha$ from $n$ includes $h$. A drifted agent may or may not have a living parent. Drift is determined by the chain's origin, not by intermediate connectivity.

The key distinction between orphaned and drifted is causal direction. Orphaning describes a disconnection in the middle of the chain — a parent is gone. Drift describes the chain's failure to originate from a human principal at all — the root is absent or non-human. An agent can be orphaned (parent gone) but not drifted (chain still traces to human). An agent can be drifted but not yet orphaned (parent alive, but chain root is non-human). The failure mode of interest arises when both conditions hold simultaneously.

\medskip
\noindent\textbf{Definition 3 (Drift Triad).}
The \emph{drift triad} is the conjunction of three conditions:
\begin{equation}
\mathrm{DriftTriad}(n) = \langle \mathrm{orphaned}(n),\ \mathrm{drifted}(n),\ \mathrm{operating}(n) \rangle.
\label{eq:rs-drift-triad}
\end{equation}
An agent enters the drift triad when all three hold: it is structurally disconnected from its parent, its chain does not reach any human anchor, and it continues executing its event loop. The third condition is essential — an orphaned, drifted agent that has terminated is not in the drift triad; it has exited. The drift triad captures the live failure state: an agent exercising consequential authority with no human in the loop.

The four possible combinations of the two binary sub-modes and operating state are shown in Table~\ref{tab:rs-drift-submodes}.

\begin{table}[h]
\centering
\begin{tabular}{c|c|c|p{5cm}}
\textbf{Orphaned} & \textbf{Drifted} & \textbf{Operating} & \textbf{Operational Meaning} \\ \hline
False & False & True  & Normal operation — chain intact, human anchor reachable \\
True  & False & True  & Orphaned but authorized — chain broken above, human anchor intact \\
False & True  & True  & Drifted but connected — chain intact, origin is non-human \\
True  & True  & True  & \textbf{Drift triad} — orphaned, drifted, operating \\
True  & True  & False & Terminated after entering drift triad (remediated) \\
True  & False & False & Normal termination following orphaning (remediated)
\end{tabular}
\caption{Drift sub-mode combinations and their operational meanings.}
\label{tab:rs-drift-submodes}
\end{table}

\medskip
\noindent\textbf{Definition 4 (Drift Heritability).}
Drift is heritable: if an agent $n$ is drifted, then any child $c$ spawned by $n$ is also drifted, provided the spawn does not undergo explicit Purpose Layer re-authorization with a human principal:
\begin{equation}
\mathrm{drifted}(n) \;\land\; \mathrm{spawned}(n, c) \;\Longrightarrow\; \mathrm{drifted}(c).
\label{eq:rs-drift-heritability}
\end{equation}
This is not a behavioral property — it is a structural consequence of the parent pointer's semantics. The child inherits the chain's drift condition through $\alpha(c) = n$, because the chain from $c$ to the root passes through $n$, which has no human anchor. One undetected drift event at depth $d$ with branching factor $b$ can produce up to $b^{d}$ drifted agents. This exponential amplification is what makes drift the central structural threat in recursive spawning systems.

% ----------------------------------------------------------------------
\subsection{Structural Conditions That Enable Drift}
\label{sec:rs-drift-conditions}

Autonomous drift arises from three structural conditions present in conventional multi-agent runtimes. None is inherently malicious; each is the consequence of architectural choices that favor operational flexibility over accountability.

\medskip
\noindent\textbf{Condition 1 — Mutable Parent Pointer (Chain Opacity).}
In conventional runtimes, $\alpha(n)$ is a mutable reference. When a parent terminates, the system may update $\alpha(n)$ to point to a new parent, or leave it pointing to the terminated node, or null it out. In any case, the original authorization trace — who authorized $n$'s provisioning and under what scope — is not preserved as a structural invariant. The runtime can report who $n$'s current parent is; it cannot guarantee that the current chain traces to the node that actually issued the original spawn warrant. Chain opacity means the delegation chain is a historical reconstruction, not a runtime-verifiable artifact.

\medskip
\noindent\textbf{Condition 2 — Unconstrained Scope Mutation.}
In conventional frameworks, an agent's operating scope $R(n)$ is re-evaluated at runtime by the agent itself, by a supervisory node, or by an external policy engine. When $n$ encounters a condition outside $R(n)$, the conventional response is scope expansion by a machine actor — the agent's authority is extended to cover the new condition and the expansion is recorded as a policy update. No human principal is consulted. The parent may not be notified. Successive expansions accumulate: each descendant inherits the expanded scope, and further expansions compound. Over $k$ spawn steps, $R(n_k)$ may bear no resemblance to $R(n_0)$.

\medskip
\noindent\textbf{Condition 3 — Silent Termination (No Child Notification).}
In conventional frameworks, agent termination is event-driven and does not require propagation to children. A parent may exit cleanly — completing its task, releasing resources — without emitting a termination event to its child nodes. The children persist in operational state, unaware that their parent has exited. They remain in the drift triad if their chain does not reach a human anchor. This is not a bug; it is the default behavior of most agent lifecycle management systems. The consequence is that child nodes can be orphaned without any structural mechanism to detect it.

These three conditions are independently necessary and jointly sufficient to produce the drift triad. A system lacking any one of them may still experience drift, but not through this particular combination of pathways.

% ----------------------------------------------------------------------
\subsection{Failure Modes}
\label{sec:rs-drift-failure-modes}

The three enabling conditions produce four concrete failure modes, each a distinct pathway to the drift triad.

\medskip
\noindent\textbf{Failure Mode 1 — Silent Orphaning.}
A parent node $\alpha(n)$ terminates without issuing a termination event to child $n$. The child $n$ remains in \texttt{ACTIVE} state, continues processing tasks, and may spawn further children — with no awareness that its parent has exited. If the chain above $\alpha(n)$ does not reach a human anchor, $n$ is also drifted. This is the Silent Orphaning pathway: the parent disappears, the child keeps working, and the chain above the parent lacks a human anchor.

\medskip
\noindent\textbf{Failure Mode 2 — Scope Creep Drift.}
A node $n$ encounters a condition $x \notin R(n)$ and — rather than escalating to its parent — expands its own scope to include $x$, recording the expansion as a policy update. The parent remains reachable; the chain topology is intact; but $R(n)$ has diverged from the human anchor's original intent. If this expansion occurs recursively at each spawn step, the effective operating scope of $n_k$ may be entirely disjoint from $R(n_0)$, even though the chain traces to a human anchor at $n_0$. The human anchor is present; the chain topology is sound; but the agent is operating outside its authorized intent. Scope creep drift is particularly insidious because the chain appears intact from a structural perspective — the drift is semantic, not topological.

\medskip
\noindent\textbf{Failure Mode 3 — Re-parenting Drift.}
A child $n$'s parent pointer $\alpha(n)$ is redirected to a new parent $p'$ by an administrative operation — for load balancing, organizational restructuring, or fault recovery. The redirection is logged as a configuration change, but the original authorization trace is not preserved. After redirection, the chain from $n$ traces to $p'$, not to the node that authorized $n$'s provisioning. If $p'$ is not itself authorized by a human principal, then $n$ has become drifted through re-parenting, even though $n$'s operational behavior has not changed. This is accountability laundering: the agent continues operating normally while the chain of authorization has been retroactively rewritten.

\medskip
\noindent\textbf{Failure Mode 4 — Ghost Authority.}
A node $n$ spawns a child $c$ without any authorized delegation chain. The spawn is initiated by a machine actor operating outside its authorized scope — a bug, a compromised Bounds Engine, or an autonomous escalation that exceeded its authority — and the child is provisioned without a valid parent pointer, without a human anchor reference, and without a Purpose Layer authorization record. The child $c$ enters operational state with a warrant that traces to no authorized chain. This failure mode requires no prior drift condition in $n$ — it can arise from a single unauthorized spawn event.

% ----------------------------------------------------------------------
\subsection{Why Depth Limits Alone Do Not Solve Drift}
\label{sec:rs-drift-depth-limits}

A common architectural response to unbounded agent spawning is a \emph{depth limit}: the system enforces a maximum chain depth $D_{\max}$, and any spawn that would produce a chain exceeding this bound is rejected. Depth limits are necessary — they prevent infinite spawn cascades and are a prerequisite for drift prevention. However, they are \emph{insufficient} to prevent autonomous drift. We prove this formally.

\medskip
\noindent\textbf{Theorem (Depth Limit Insufficiency).}
Let $D_{\max}$ be a system-enforced maximum chain depth. Then there exists at least one execution path in which a chain of depth $\leq D_{\max}$ produces a drifted agent, and no choice of $D_{\max}$ prevents this without additional structural constraints on scope refinement.

\medskip
\noindent\textbf{Proof.} Consider a spawn chain $C = \langle n_0, n_1, \ldots, n_k \rangle$ where $n_0 = h \in \mathcal{H}$ is a human principal and $k < D_{\max}$. The chain satisfies the depth constraint by construction: $\mathrm{depth}(n_k) = k \leq D_{\max}$. Apply Failure Mode 2 (Scope Creep Drift): each node $n_i$ for $1 \leq i \leq k$ successively expands its operating scope $R(n_i)$ beyond $R(n_{i-1})$ without Purpose Layer authorization, recording each expansion as a policy update. After $k$ successive expansions, $R(n_k)$ is entirely disjoint from $R(n_0)$. The chain depth is $k \leq D_{\max}$; the chain has a human anchor at $n_0$; yet $n_k$ is drifted, because its operating scope $R(n_k)$ is not a descendant refinement of $h(n_0)$'s intent. No choice of $D_{\max}$ prevents this, because the depth constraint operates on the length of the chain, not on the scope relationship between successive nodes. $\blacksquare$

The drift prevention guarantee requires the conjunction of two independent constraints:
\begin{equation}
\mathrm{Drift\ Prevention} \;\Longleftrightarrow\; \bigl( \mathrm{depth}(n) \leq D_{\max} \bigr)\ \land\ \bigl( R(n_{i+1}) \subseteq R(n_i)\ \forall\, i \bigr).
\label{eq:rs-drift-prevention-constraint}
\end{equation}
The first conjunct is the depth constraint, enforced by the Bounds Engine. The second conjunct is the scope refinement constraint: each child node's operating scope must be a subset of its parent's, ensuring the chain is a monotonically narrowing refinement of the human anchor's intent. Depth limits enforce only the first conjunct. A system enforcing $D_{\max}$ without the scope refinement constraint is immune to unbounded spawn cascades but remains fully vulnerable to scope creep drift.

The deeper structural reason depth limits are insufficient is that they address the wrong dimension of the accountability problem. Depth limits constrain chain length — a structural property of the spawn graph. Drift, however, is fundamentally about authorization scope: whether each node's operating envelope is a legitimate descendant refinement of the human anchor's intent. A drifted agent at depth 1 with unbounded scope expansion authority can produce consequences that exceed the human anchor's intent as severely as a drifted agent at depth $D_{\max} - 1$.

% ----------------------------------------------------------------------
\subsection{Formal Model}
\label{sec:rs-drift-model}

We define a compact formal model capturing the structural elements essential to drift analysis. This model is referenced in subsequent sections and used in the correctness proofs.

\medskip
\noindent\textbf{System Tuple.}
The system $\mathcal{S}$ is a 5-tuple:
\begin{equation}
\mathcal{S} = \langle \mathcal{N},\ \alpha,\ R,\ D_{\max},\ \mathcal{L} \rangle.
\label{eq:rs-system-model}
\end{equation}
$\mathcal{N}$ is the set of all agent nodes. $\alpha: \mathcal{N} \rightarrow \mathcal{N} \cup \{\bot\}$ is the parent pointer function, where $\alpha(n) = \bot$ for root human principals. $R: \mathcal{N} \rightarrow 2^{\mathcal{A}}$ maps each node to its operating scope as a set of authorized actions. $D_{\max}$ is the system-enforced maximum chain depth. $\mathcal{L}$ is the forensic ledger recording all authorized spawn events.

\medskip
\noindent\textbf{Chain Depth.}
The depth of a node $n$ is:
\begin{equation}
\mathrm{depth}(n) =
\begin{cases}
0 & \text{if } \alpha(n) = \bot \quad (\text{root}) \\
1 + \mathrm{depth}(\alpha(n)) & \text{otherwise}.
\end{cases}
\label{eq:rs-depth}
\end{equation}

\medskip
\noindent\textbf{Human Anchor Reachability.}
The predicate $\mathrm{chain-reaches}(n, h)$ returns true if the transitive closure of $\alpha$ from $n$ includes the human anchor $h$:
\begin{equation}
\mathrm{chain-reaches}(n, h) \;\Longleftrightarrow\; \exists\, \text{path } \langle n = n_k, n_{k-1}, \ldots, n_0 = h \rangle \text{ s.t. } \alpha(n_i) = n_{i-1}.
\label{eq:rs-chain-reaches}
\end{equation}

\medskip
\noindent\textbf{Drift Condition.}
An agent $n$ is in the drift triad if and only if:
\begin{equation}
\bigl( \alpha(n) = \bot \ \lor\ \mathrm{unreachable}(\alpha(n)) \bigr)\ \land\ \bigl( \neg\exists\, h \in \mathcal{H} : \mathrm{chain-reaches}(n, h) \bigr)\ \land\ \mathrm{operating}(n) = \mathrm{true}.
\label{eq:rs-drift-condition}
\end{equation}
The first conjunct establishes orphaning; the second establishes that no human anchor is reachable; the third establishes continued operation.

\medskip
\noindent\textbf{Scope Refinement Invariant (RSP).}
Under the Recursive Spawning Protocol, every spawn from $n_i$ to $n_{i+1}$ must satisfy:
\begin{equation}
R(n_{i+1}) \subseteq R(n_i) \quad \text{and} \quad \mathrm{depth}(n_{i+1}) \leq D_{\max}.
\label{eq:rs-scope-refinement}
\end{equation}
Both are enforced by the Fact Layer pre-commit hook at provisioning time, making them structural invariants rather than policy conventions.

This model provides the vocabulary for the Recursive Spawning Protocol specification in Section~\ref{sec:rs-protocol}, which modifies $\mathcal{S}$ by replacing the mutable parent pointer and scope functions with their immutable counterparts, and by adding the pre-commit hook and forensic ledger that enforce Equation~\ref{eq:rs-scope-refinement} as architectural invariants.